\subsection{Environments}
\label{sec:environments}

Pat's type system makes use of \emph{environments} for keeping track of types of variables while type checking terms.
The implementation of environments depends on how variables are represented.
For traditional strings as variables, it is common to use a map, mapping variables to types \cite{pierceTypesProgrammingLanguages2002}.
In the case of de Bruijn indices, an environment is typically represented as a list of types where a variable indicates the position of its type in the list \cite{pierceTypesProgrammingLanguages2002,wadlerProgrammingLanguageFoundations2022a}.
%
For Pat, the representation of an environment as a simple list of types is
unfortunately not enough, due to its parallel fragment and use of
quasilinearity.  To illustrate the problem, consider the type derivation of a
mailbox calculus term on the left-hand side:

{
    \small
\begin{mathpar}
    \inferrule*
    { A_1, \emptyset \vdash \Var{0} \\ {A_2, \emptyset \vdash \Var{1}} }
    {A_2, A_1, \emptyset \vdash \Var{0} \parallel \Var{1}}

    \inferrule*
     { \bot, A_1, \emptyset \vdash \Var{0} \\
       A_2, \bot, \emptyset \vdash \Var{1}}
     {A_2, A_1, \emptyset \vdash \Var{0} \parallel \Var{1}}
\end{mathpar}
}

We assume variable $\Var{0}$ has type $A_1$ and variable $\Var{1}$ has type $A_2$.
The term expresses that there are two processes running in parallel, each with only a single variable.
Additionally, the current environment consists of $A_2, A_1, \emptyset$.
By the typing rules of the mailbox calculus (read bottom-up), given by de'Liguoro and Padovani \cite{deliguoroMailboxTypesUnordered2018}, when typing parallel processes, the environment needs to be split into two.
Assuming variable $\Var{0}$ has indeed type $A_1$, the left subtree works without problems.
Looking up the type at position $\Var{0}$ in environment $A_1, \emptyset$ yields $A_1$.
However, the right subtree poses a problem.
There is no type at position $\Var{1}$ in environment $A_2, \emptyset$, as it only contains a single element.

One way of solving this issue is the inclusion of entries in the environment representing the absence of a type at a certain position.
The example on the right-hand side adds an additional $\bot$-entry, representing an empty position.
%
With this approach, whenever an environment is split, the size of the two new environments is equal to the size of the original one.
%
The \emph{dblib} library we use for managing de Bruijn indices also includes an implementation of environments based on this exact approach.
It provides the needed definitions and operations, such as lookup and insertion, together with proofs about properties on environments.

\begin{figure}[t]
    \small
  \textbf{Environment subtyping} (\rocqlink{https://edgardschi.github.io/mailbox-types-rocq/MailboxTypes.Environment.html\#EnvironmentSubtype}) \hfill \fbox{$\Gamma_1 \subtype \Gamma_2$}
  \begin{center}
        \begin{prooftree}
            \infer0{\Gamma \subtype \Gamma}
        \end{prooftree}
      \;
        % \quad
        \begin{prooftree}
            \hypo{\Gamma_1 \subtype \Gamma_2}
            \hypo{\Gamma_2 \subtype \Gamma_3}
            \infer2{\Gamma_1 \subtype \Gamma_3}
        \end{prooftree}
      \;
        % \quad
        \begin{prooftree}
            \hypo{\Gamma_1 \subtype \Gamma_2}
            \infer1{\bot, \Gamma_1 \subtype \bot, \Gamma_2}
        \end{prooftree}
      \;
        % \quad
        \begin{prooftree}
            \hypo{\Gamma_1 \subtype \Gamma_2}
            \hypo{\unr{A}}
            \infer2{A, \Gamma_1 \subtype \bot, \Gamma_2}
        \end{prooftree}
      \;
        % \quad
        \begin{prooftree}
            \hypo{\Gamma_1 \subtype \Gamma_2}
            \hypo{A_1 \subtype A_2}
            \infer2{A_1, \Gamma_1 \subtype A_2, \Gamma_2}
        \end{prooftree}
    \end{center}
  \textbf{Environment combination} (\rocqlink{https://edgardschi.github.io/mailbox-types-rocq/MailboxTypes.Environment.html\#EnvironmentCombination}) \hfill \fbox{$\mbtucomb{\Gamma_1}{\Gamma_2}{\Gamma_3}$}
  \begin{center}
        \begin{prooftree}
            \infer0{\mbtucomb{\emptyset}{\emptyset}{\emptyset}}
        \end{prooftree}
        \quad
        \begin{prooftree}
            \hypo{\mbtucomb{\Gamma_1}{\Gamma_2}{\Gamma_3}}
            \infer1{\mbtucomb{\bot, \Gamma_1}{\bot, \Gamma_2}{\bot, \Gamma_3}}
        \end{prooftree}
        \quad
        \begin{prooftree}
            \hypo{\mbtucomb{\Gamma_1}{\Gamma_2}{\Gamma_3}}
            \infer1{\mbtucomb{A, \Gamma_1}{\bot, \Gamma_2}{A, \Gamma_3}}
        \end{prooftree}
        \quad
        \begin{prooftree}
            \hypo{\mbtucomb{\Gamma_1}{\Gamma_2}{\Gamma_3}}
            \infer1{\mbtucomb{\bot, \Gamma_1}{A, \Gamma_2}{A, \Gamma_3}}
        \end{prooftree}
    \end{center}
    \begin{center}
        \begin{prooftree}
            \hypo{\mbtucomb{A_1}{A_2}{A_3}}
            \hypo{\mbtucomb{\Gamma_1}{\Gamma_2}{\Gamma_3}}
            \infer2{\mbtucomb{A_1, \Gamma_1}{A_2, \Gamma_2}{A_3, \Gamma_3}}
        \end{prooftree}
    \end{center}
  \textbf{Disjoint environment combination} (\rocqlink{https://edgardschi.github.io/mailbox-types-rocq/MailboxTypes.Environment.html\#EnvironmentDisjointCombination}) \hfill \fbox{$\disenvcomb{\Gamma_1}{\Gamma_2}{\Gamma_3}$}
  \begin{center}
        \begin{prooftree}
            \infer0{\disenvcomb{\emptyset}{\emptyset}{\emptyset}}
        \end{prooftree}
        \quad
        \begin{prooftree}
            \hypo{\disenvcomb{\Gamma_1}{\Gamma_2}{\Gamma_3}}
            \infer1{\disenvcomb{\bot, \Gamma_1}{\bot, \Gamma_2}{\bot, \Gamma_3}}
        \end{prooftree}
        \quad
        \begin{prooftree}
            \hypo{\disenvcomb{\Gamma_1}{\Gamma_2}{\Gamma_3}}
            \infer1{\disenvcomb{A, \Gamma_1}{\bot, \Gamma_2}{A, \Gamma_3}}
        \end{prooftree}
        \quad
        \begin{prooftree}
            \hypo{\disenvcomb{\Gamma_1}{\Gamma_2}{\Gamma_3}}
            \infer1{\disenvcomb{\bot, \Gamma_1}{A, \Gamma_2}{A, \Gamma_3}}
        \end{prooftree}
    \end{center}
    \begin{center}
        \begin{prooftree}
            \hypo{\disenvcomb{\Gamma_1}{\Gamma_2}{\Gamma_3}}
            \infer1{\disenvcomb{C, \Gamma_1}{C, \Gamma_2}{C, \Gamma_3}}
        \end{prooftree}
    \end{center}

    \caption{The definitions of relations on environments.}
    \label{fig:environment-relations}
\end{figure}

% \begin{linkdef}{https://edgardschi.github.io/mailbox-types-rocq/MailboxTypes.Environment.html\#Env}[Environment]
%   \todo{ES}{Move this into syntax figure perhaps}
%   An \emph{environment} $\Gamma$ is defined by the following grammar:
%
%   \begin{center}
%     \begin{grammar}
%       \firstcasesubtil{$\Gamma$}{\emptyset \hspace{-.2cm}\gralt\hspace{-.2cm} \bot, \Gamma \hspace{-.2cm}\gralt\hspace{-.2cm} A, \Gamma}{}
%     \end{grammar}
%   \end{center}
% \end{linkdef}

% In Rocq, an environment is a list with elements of type \coqe{option TUsage}, i.e.\ elements can either be \coqe{None} or \coqe{Some A} for some usage-annotated type \coqe{A}.
% Note how the syntax forces an order onto the elements of an environment.
% Reading from right to left, the first element (after $\emptyset$) corresponds to position $0$, while for any element at position $n$ its neighbor to the left has position $n+1$.
%
% To reason about elements at arbitrary positions in an environment, as well as inserting them at a given position, \emph{dblib} includes an \emph{insert} operation.
%
% \begin{linkdef}{https://edgardschi.github.io/mailbox-types-rocq/MailboxTypes.Dblib.Environments.html\#raw_insert}[Environment Insert]
%   \emph{Inserting} a new variable $x$ associated with element $e$ in environment $\Gamma$ (written $\envInsert{x}{e}{\Gamma}$) is defined by the following equations:
%   \begin{alignat*}{3}
%     &\envInsert{0}{e}{\Gamma} &&:= e, \Gamma\\
%     &\envInsert{x+1}{e}{(f, \Gamma)} &&:= f, (\envInsert{x}{e}{\Gamma})\\
%     &\envInsert{x+1}{e}{\emptyset} &&:= \bot, (\envInsert{x}{e}{\emptyset})\\
%   \end{alignat*}
% \end{linkdef}

% Essentially, what the operation $\envInsert{x}{e}{\Gamma}$ does is to move $x$ positions into $\Gamma$ and insert $e$.
% If the empty environment $\emptyset$ is encountered while traversing, a new $\bot$ is added each time, building up the environment.
%
% Additionally, \emph{dblib} provides an implementation for the \emph{lookup} function, taking as arguments a variable and environment and returning either $\bot$ or a type.
%
% \begin{linkdef}{https://edgardschi.github.io/mailbox-types-rocq/MailboxTypes.Dblib.Environments.html\#lookup}[Environment Lookup]
%   \emph{Looking up} the type for a variable $x$ in environment $\Gamma$ is defined as follows:
%   \begin{alignat*}{3}
%     &lookup \ x \ \emptyset &&:= \bot\\
%     &lookup \ 0 \ (e,\Gamma) &&:= e\\
%     &lookup \ (x + 1) \ (e,\Gamma) &&:= lookup \ x \ \Gamma
%   \end{alignat*}
% \end{linkdef}
%

Fowler et al. \cite{fowlerSpecialDeliveryProgramming2023} describe several relations and operations on environments, defined in Figure \ref{fig:environment-relations}.
Our definitions are essentially equal to the original ones, with modifications to incorporate $\bot$-entries and some corrections that are further discussed at the end of this section.
In each relation we add a rule for the case where the first entry in both environments is $\bot$.
Similar to partial operations on types, we turn partial operations on environments into relations.

% Having defined the basics for environments, we continue with relations and operations on environments, as described by Fowler et al. \cite{fowlerSpecialDeliveryProgramming2023}.
% The following definitions are based on Pat's original definitions, with a few modifications for our environment representation.
%
% First, we extend subtyping to environments, which also includes a notion of \emph{weakening}.

% \begin{linkdef}{https://edgardschi.github.io/mailbox-types-rocq/MailboxTypes.Environment.html\#EnvironmentSubtype}[Environment Subtyping]
%   \emph{Environment subtyping} (written $\Gamma_1 \subtype \Gamma_2$) is defined by the following rules:
%   \begin{center}
%         \begin{prooftree}
%             \infer0{\Gamma \subtype \Gamma}
%         \end{prooftree}
%         \qquad
%         \begin{prooftree}
%             \hypo{\Gamma_1 \subtype \Gamma_2}
%             \hypo{\Gamma_2 \subtype \Gamma_3}
%             \infer2{\Gamma_1 \subtype \Gamma_3}
%         \end{prooftree}
%         \qquad
%         \begin{prooftree}
%             \hypo{\Gamma_1 \subtype \Gamma_2}
%             \infer1{\bot, \Gamma_1 \subtype \bot, \Gamma_2}
%         \end{prooftree}
%         \qquad
%         \begin{prooftree}
%             \hypo{\Gamma_1 \subtype \Gamma_2}
%             \hypo{\unr{A}}
%             \infer2{A, \Gamma_1 \subtype \bot, \Gamma_2}
%         \end{prooftree}
%     \end{center}
%     \begin{center}
%         \begin{prooftree}
%             \hypo{\Gamma_1 \subtype \Gamma_2}
%             \hypo{A_1 \subtype A_2}
%             \infer2{A_1, \Gamma_1 \subtype A_2, \Gamma_2}
%         \end{prooftree}
%     \end{center}
% \end{linkdef}

% The definition for subtyping of environments includes an additional rule specifically for our representation of environments, where $\bot$ entries at the same position in both environments are skipped.
% Otherwise, the rules follow the definition of Fowler et al., with some corrections that are further discussed in Chapter \ref{sec:rocq-mechanization}.
% It follows directly from these rules that environment subtyping is reflexive and transitive.

% In addition to environment subtyping, we also include the notion of \emph{strict environment subtyping}, where the domains of both environments must be equal.
% Meaning that, in our representation, there is either a type or $\bot$ in both environments at the same positions.
%
% \begin{linkdef}{https://edgardschi.github.io/mailbox-types-rocq/MailboxTypes.Environment.html\#EnvironmentSubtypeStrict}[Strict Environment Subtyping]
%   \emph{Strict environment subtyping} (written $\Gamma_1 \strictsubtype \Gamma_2$) is defined by the following rules:
%   \begin{center}
%         \begin{prooftree}
%             \infer0{\Gamma \strictsubtype \Gamma}
%         \end{prooftree}
%         \qquad
%         \begin{prooftree}
%             \hypo{\Gamma_1 \strictsubtype \Gamma_2}
%             \hypo{\Gamma_2 \strictsubtype \Gamma_3}
%             \infer2{\Gamma_1 \strictsubtype \Gamma_3}
%         \end{prooftree}
%         \qquad
%         \begin{prooftree}
%             \hypo{\Gamma_1 \strictsubtype \Gamma_2}
%             \infer1{\bot, \Gamma_1 \strictsubtype \bot, \Gamma_2}
%         \end{prooftree}
%         \qquad
%         \begin{prooftree}
%             \hypo{\Gamma_1 \strictsubtype \Gamma_2}
%             \hypo{A_1 \strictsubtype A_2}
%             \infer2{A_1, \Gamma_1 \strictsubtype A_2, \Gamma_2}
%         \end{prooftree}
%     \end{center}
% \end{linkdef}


% \begin{linkdef}{https://edgardschi.github.io/mailbox-types-rocq/MailboxTypes.Environment.html\#EnvironmentCombination}[Environment Combination]
%   \emph{Environment combination} (written $\mbtucomb{\Gamma_1}{\Gamma_2}{\Gamma_3}$) is defined as a relation between three environments, by the following rules:
%   \begin{center}
%         \begin{prooftree}
%             \infer0{\mbtucomb{\emptyset}{\emptyset}{\emptyset}}
%         \end{prooftree}
%         \qquad
%         \begin{prooftree}
%             \hypo{\mbtucomb{\Gamma_1}{\Gamma_2}{\Gamma_3}}
%             \infer1{\mbtucomb{\bot, \Gamma_1}{\bot, \Gamma_2}{\bot, \Gamma_3}}
%         \end{prooftree}
%         \qquad
%         \begin{prooftree}
%             \hypo{\mbtucomb{\Gamma_1}{\Gamma_2}{\Gamma_3}}
%             \infer1{\mbtucomb{A, \Gamma_1}{\bot, \Gamma_2}{A, \Gamma_3}}
%         \end{prooftree}
%     \end{center}
%     \begin{center}
%         \begin{prooftree}
%             \hypo{\mbtucomb{\Gamma_1}{\Gamma_2}{\Gamma_3}}
%             \infer1{\mbtucomb{\bot, \Gamma_1}{A, \Gamma_2}{A, \Gamma_3}}
%         \end{prooftree}
%         \qquad
%         \begin{prooftree}
%             \hypo{\mbtucomb{A_1}{A_2}{A_3}}
%             \hypo{\mbtucomb{\Gamma_1}{\Gamma_2}{\Gamma_3}}
%             \infer2{\mbtucomb{A_1, \Gamma_1}{A_2, \Gamma_2}{A_3, \Gamma_3}}
%         \end{prooftree}
%     \end{center}
% \end{linkdef}

Pat requires subtyping on \emph{environments} rather than just types in order to
rewrite types so that they can be used by the type combination operator.
Sequential environment combination $\envcomb{\Gamma_1}{\Gamma_2}{\Gamma}$ makes
use of type combination and is used when type checking let-expressions. This is
the main rule used to ensure that subsequent uses of a mailbox variable
``balance out''.
Additionally, we use \emph{disjoint environment combination} to reason about
expressions where subterms may not share mailbox-typed variables (e.g.\ a
variable may not be used as the target of a message send and also be contained
in its payload). Environments combined using $+$ are only allowed to share
variables of base types and be disjoint otherwise.

As our definition of environments includes $\bot$ entries, it is possible for an environment to be empty without being $\emptyset$, necessitating a new notion of emptiness of environments $\EmptyEnv{\Gamma}$ (\rocqlink{https://edgardschi.github.io/mailbox-types-rocq/MailboxTypes.Environment.html\#EmptyEnv})
that holds when $\Gamma$ only contains $\bot$-entries.

\paragraph{Mechanization.}
Formalizing environments is, by far, the most complex part of this mechanization.
The library \emph{dblib} reduces some of the complexity around environments by proving operations and properties on them.
However, reasoning about environment combinations and environment subtyping still is far from trivial.
Independent of the representation of environments that were considered, environment splitting and subtyping require dozens of technical lemmas.
While the proofs of these properties follow the same general approach, they are lengthy and tedious to perform.
%
For example, for the relation $\envcomb{\Gamma_1}{\Gamma_2}{\Gamma_3}$ we have proved that
there exists an environment that is the environment combination of both
$\Gamma_1$ and $\Gamma_2$ with the entry at the same position replaced by $\bot$ in both environments (\rocqlink{https://edgardschi.github.io/mailbox-types-rocq/MailboxTypes.Environment.html\#EnvironmentCombination_tes}).
Similar lemmas had to be proven for environment subtyping and disjoint combination.
While mechanizing subtyping for environments according to the definition of Pat \cite{fowlerSpecialDeliveryProgramming2023},
we noticed that the relation was missing an inference rule, making environment subtyping not transitive.
By explicitly including the transitivity rule we were able to prove desired properties of environment subtyping.
This issue was subsequently fixed in a revised version \cite{fowlerSpecialDeliveryProgramming2025}.
% Not the journal version yet -- still under review :)

